\uuid{d175}
\exo7id{551}
\titre{exo7 551}
\auteur{cousquer}
\organisation{exo7}
\datecreate{2003-10-01}
\isIndication{false}
\isCorrection{false}
\chapitre{Suite}
\sousChapitre{Suite définie par une relation de récurrence}
\cpge{mpsi-mp2i.nombres-reels-suites.c, pcsi.nombres-reels-suites.b, ptsi.nombres-reels-suites.b, mpsi-mp2i.nombres-reels-suites.d, pcsi.nombres-reels-suites.c, ptsi.nombres-reels-suites.c, mpsi-mp2i.nombres-reels-suites.e, pcsi.nombres-reels-suites.d, ptsi.nombres-reels-suites.d, mpsi-mp2i.nombres-reels-suites.i, pcsi.nombres-reels-suites.g, ptsi.nombres-reels-suites.g, bcpst1.suites-reelles, bcpst1.limites-continuite.a}
\module{Analyse}
\niveau{L1}
\difficulte{}

\contenu{
\texte{
Soient $a$ et $b$ deux réels strictement positifs~; on définit
une suite $(u_n)$ par~:
$$ u_0 \ge0 \quad{\rm et} \quad u_{n+1}= \sqrt{au_n+b}$$
}
\begin{enumerate}
    \item \question{Montrer qu'il existe une valeur de $u_0$ pour laquelle cette suite
est stationnaire.}
    \item \question{Montrer que si $u_0$ est distinct de cette valeur, $(u_n)$ est monotone et
bornée. Trouver $\lim_{n \to\infty} u_n $.}
\end{enumerate}
}
